% AN OUTPUT ROUTINE RUNS IN INTERNAL VERTICAL MODE, AND ENDS WITH A BRACE.
%
% \tracingcommands names the mode where it changes, so an \output of
%
%   {\count9=1 \shipout\box255 \count9=2 }
%
% draws
%
%   {vertical mode: \penalty}          <- what set the page builder going
%   {internal vertical mode: \count}   <- the routine's first command
%   {\shipout}
%   {\count}
%   {end-group character }}            <- the routine's closing brace
%   {vertical mode: \message}          <- back where the page builder was
%
% Two things there. The mode is INTERNAL vertical, not the vertical mode
% the page builder was in -- the routine builds a list of its own. And
% the routine is closed by a RIGHT BRACE that the reference obeys and
% traces like any other, one per output routine.
%
% hstex ran the routine in the mode that stood and closed its group
% without a token, so it drew neither line; and because the mode had not
% changed, the next command after the routine did not name its mode
% either. That third line is the one that showed up in trip -- fifty-odd
% missing mode prefixes downstream of a routine that had run.
\tracingonline=1
\output={\count9=1 \shipout\box255 \count9=2 }
\vsize=10pt \maxdeadcycles=3
\tracingcommands=2
X\par\penalty-10000
\message{[done]}
\tracingcommands=0
\end
